\section{Natura della Luce}

\begin{frame}{Cos'è la luce}
    \begin{itemize}
        \item La luce è una forma di enegia che si propaga nello spazio
        \item È responsabile della visione e della percezione dei colori e permette il trasferimento di energia e informazioni
        \item Per secoli si è discusso sulla natura della luce e si presentano due modelli:
        \begin{itemize}
            \item modello corpuscolare (particelle)
            \item modello ondulatorio (onde)
        \end{itemize}
    \end{itemize}
\end{frame}

\newcommand{\currentlink}{} % URL di base

%=======================================================================

\begin{frame}{Dualismo onda-particella}
    \begin{itemize}
        \item La luce presenta una doppia natura:
        \begin{itemize}
            \item si comporta da particella (fotone) quando interagisce con la materia
            \item si comporta da onda quando si propaga nello spazio
        \end{itemize}
        \item Non è una contraddizione: sono due descrizioni complementari dello stesso fenomeno
    \end{itemize}
\end{frame}


%=======================================================================

\begin{frame}{L'Onda Elettromagnetica}
    \begin{itemize}
        \item La luce può essere descritta come un'onda precisamente come un'onda elettromagnetica
        \item Onda trasversale: campo elettrico $\vec{E}$ e campo magnetico $\vec{B}$ oscillano perpendicolarmente alla direzione di propagazione
        \item Non richiede un mezzo materiale: si propaga nel vuoto
        \item Relazione fondamentale: 
        $v = \lambda f$ \\
        dove $\lambda$ = lunghezza d'onda, $f$ = frequenza, $v$ = velocità di propagazione
        \item $\lambda$ e $f$ sono inversamente proporzionali a velocità fissata
    \end{itemize}
\end{frame}




